\section{Introduction}

\begin{sloppypar}
This paper presents \lafc, a benchmark for learned cache eviction built from real-trace-derived cache decisions and computed finite-horizon counterfactual labels. The contribution is not a new eviction policy. It is a reusable benchmark surface for caching research that exposes candidate-level supervision, decision-level summaries, a shipped pairwise sample, and release-oriented metadata and validation. The motivating gap is now familiar in learned systems work: policy papers often compare methods under different traces, different preprocessing conventions, different supervision targets, and different release practices, which makes cross-paper empirical claims harder to interpret than the headline numbers suggest.

\lafc\ is designed to address that gap by making the supervised unit of comparison explicit. Each canonical candidate row corresponds to one candidate victim at one full-cache eviction decision, with decision metadata, candidate features, and a continuation-policy-specific finite-horizon counterfactual label. The release stores this label as \texttt{y\_loss}, defines \texttt{y\_value} as $-\texttt{y\_loss}$, and treats both as properties of a candidate within a decision rather than as properties of a fully optimized future policy. From the same underlying release, the benchmark also exposes a decision view and a pairwise sample so that value prediction, within-decision best-candidate selection, pairwise preference prediction, and label-distribution characterization can be studied under a shared artifact. In that sense, \lafc\ is not merely a raw trace repository either: it sits above the trace layer and packages a decision-aligned supervision layer for benchmark reuse.

The evaluated open-release snapshot is already large enough to make that benchmark framing substantive. Table~\ref{tab:dataset-scale} reports manifest-backed release scale of 277,995,072 candidate rows, 2,363,286 decision-view rows, and a shipped pairwise sample of 1,000,000 rows across five selected trace families. Table~\ref{tab:decision-breakdown} summarizes the decision view by family, capacity, horizon, and split. Table~\ref{tab:candidate-count-stats} shows that decision complexity is nontrivial even before any model training: in the decision view, candidate counts range from 32 to 256, with median 64 and mean 117.63. Table~\ref{tab:pairwise-sample-stats} reports coverage of 118,635 unique decisions in the pairwise sample. These numbers scope the benchmark, and the evaluated release package passed end-to-end validation. The paper complements those release descriptors with row-level \texttt{y\_loss}/\texttt{y\_value} summaries and lightweight baselines for value regression, best-candidate selection, and pairwise preference prediction.
\end{sloppypar}

The release also makes a deliberate boundary between raw data provenance and benchmark packaging. \lafc\ does not claim authorship of the underlying traces, and the open subset is intentionally narrower than the broader generated corpus. Instead, the benchmark contribution is the releaseable supervision artifact derived from those traces. It includes the schema, the task views, the manifest, the checksums, and the validation machinery needed for reuse. That separation matters scientifically. The labels are finite-horizon and continuation-policy-specific decision targets rather than offline-optimal targets. It also matters operationally because redistribution constraints affect what can responsibly appear in the open subset.

To our knowledge, existing public cache-replacement artifacts do not provide a reusable candidate-level counterfactual supervision table at this scale. The novelty claim is therefore about the released supervision artifact, schema, validation surface, and benchmark tasks rather than about a new eviction policy.

\begin{sloppypar}
The paper therefore makes a benchmark claim rather than a deployment claim. First, learned cache eviction can be studied through decision-aligned candidate supervision instead of only through end-policy comparisons. Second, one release can support multiple related tasks through candidate, decision, and pairwise views. Third, release engineering and validation are part of the scientific contribution because they determine whether later work can reproduce or extend the benchmark. Fourth, the paper focuses on benchmark construction, release structure, characterization, and simple baselines without overstating downstream systems impact or claiming model coverage beyond the evidence reported here. The baselines are deliberately modest in that sense: linear value regression remains weak out of sample, the induced best-candidate selector is useful as a reproducible benchmark floor rather than as evidence of production-quality policy behavior, and the new pairwise results show decision-relevant signal in released candidate features without converting the paper into a full-policy study.
\end{sloppypar}

The rest of the paper proceeds as follows. Section~\ref{sec:background} reviews the learned cache-eviction setting and motivates decision-aligned supervision. Section~\ref{sec:benchmark-design} defines the benchmark scope and release-governance boundary. Sections~\ref{sec:label-generation} and~\ref{sec:schema-views} describe the label definition, schema, and derived views. Section~\ref{sec:release-validation} explains release construction and validation. Sections~\ref{sec:tasks-baselines} and~\ref{sec:characterization} define the benchmark tasks and the current characterization surface. Section~\ref{sec:limitations} covers limitations, ethics, and reuse constraints, followed by related work and conclusion.
